\documentclass[10pt]{article}
\usepackage[T1]{fontenc}
\usepackage[utf8]{inputenc}
\usepackage[french]{babel}
\usepackage{amsmath,amssymb,mathrsfs,mathpazo}
\usepackage{geometry}
\geometry{a4paper,margin=2cm}
\usepackage{graphicx,array,tabularx,tabularray,colortbl,multirow,multicol,booktabs,enumitem,pifont,ulem,fancybox,diagbox,cancel}
\usepackage{pstricks,pstricks-add,pst-plot,pst-tree,pst-node,pst-text,pst-func,pst-math,pst-bspline,pst-3dplot,pst-eucl,pst-solides3d}
\let\pscalculate\undefined
\usepackage{pst-bezier}
\usepackage{tikz,pgfplots,tkz-tab}
\usetikzlibrary{arrows,arrows.meta,patterns,calc,positioning,trees,decorations.markings,shapes.geometric,shapes.misc,shapes.arrows,intersections,decorations.pathmorphing,decorations.pathreplacing,backgrounds,shadows,babel}
\pgfplotsset{compat=1.18}
\usepgfplotslibrary{fillbetween}
\usepackage[np]{numprint}
\usepackage{listings,makecell,pas-tableur,twoopt}
\let\barre\undefined
\usepackage[pstricks]{tablvar}
\definecolor{darkgreen}{RGB}{0,100,0}
\providecommand{\red}{\color{red}}
\providecommand{\blue}{\color{blue}}
\providecommand{\green}{\color{green}}
\providecommand{\black}{\color{black}}
\providecommand{\white}{\color{white}}
\providecommand{\gray}{\color{gray}}
\providecommand{\cyan}{\color{cyan}}
\providecommand{\magenta}{\color{magenta}}
\providecommand{\yellow}{\color{yellow}}
\providecommand{\violet}{\color{violet}}
\providecommand{\orange}{\color{orange}}
\providecommand{\euro}{\texteuro}
\providecommand{\EUR}{\texteuro}
\providecommand{\ieme}{\textsuperscript{e}}
\providecommand{\degres}{\ensuremath{^\circ}}
\providecommand{\degree}{\ensuremath{^\circ}}
\providecommand{\vv}[1]{\overrightarrow{#1}}
\providecommand{\Coord}[1]{#1}
\providecommand{\decosix}{}
\providecommand{\starredbullet}{\ensuremath{\bullet}}
\pagestyle{empty}
\setlength{\parindent}{0pt}
\usepackage{adjustbox}\DeclareUnicodeCharacter{2236}{:}\DeclareUnicodeCharacter{2190}{\ensuremath{\leftarrow}}
\providecommand{\R}{\mathbb{R}}
\providecommand{\N}{\mathbb{N}}
\providecommand{\Z}{\mathbb{Z}}
\providecommand{\Q}{\mathbb{Q}}
\providecommand{\D}{\mathbb{D}}
\providecommand{\C}{\mathbb{C}}
\providecommand{\e}{{\mathrm{e}}}
\providecommand{\vect}[1]{\overrightarrow{#1}}
\providecommand{\vectt}[1]{\overrightarrow{\mathrm{#1}}}
\providecommand{\vvt}[1]{\overrightarrow{\mathrm{#1}}}
\providecommand{\np}[2][]{{#2~#1}}
\providecommand{\vv}[1]{{\overrightarrow{#1}}}
\renewcommand{\d}{{\mathrm{d}}}
\providecommand{\degres}{{^\circ}}
\providecommand{\degre}{{^\circ}}
\providecommand{\blue}{\color{blue}}
\providecommand{\red}{\color{red}}
\providecommand{\enskip}{\quad}
\providecommand{\llbracket}{\lbrack\!\lbrack}
\providecommand{\rrbracket}{\rbrack\!\rbrack}
\providecommand{\euro}{{\text{€}}}
\providecommand{\no}{{\text{n°}}}
\providecommand{\interval}[2]{{[#1;#2]}}
\providecommand{\intervFO}[2]{{[#1;#2[}}
\providecommand{\dotfill}{\ldots}

\providecommand{\ds}{\displaystyle}
\providecommand{\textding}[1]{#1}
\providecommand{\Oij}{\(\mathrm{O};\vec{\imath},\vec{\jmath}\)}
\providecommand{\Oijk}{\(\mathrm{O};\vec{\imath},\vec{\jmath},\vec{k}\)}
\providecommand{\Ouv}{\(\mathrm{O};\vec{u},\vec{v}\)}
\providecommand{\intervFF}[2]{\ensuremath{[#1;#2]}}
\renewcommand{\euro}{\ifmmode\text{\texteuro}\else\texteuro\fi}
\DeclareUnicodeCharacter{00B0}{\ifmmode^{\circ}\else\textdegree\fi}
\DeclareUnicodeCharacter{20AC}{\ifmmode\text{\texteuro}\else\texteuro\fi}
% Énoncés : Éducation nationale. Transcription et corrigés : APMEP, Denis Vergès et contributeurs. Découpage et classement : CoopMaths, Sébastien Lozano, Éric Elter et contributeurs. Adaptation de la présentation : Maths Explorer.
% CC BY-SA : https://coopmaths.fr/www/about/
% https://coopmaths.fr/alea/static/e3c/2020/tex/e3c_2020_00_sujet61_4.tex
\begin{document}
\noindent\textit{APMEP / CoopMaths. CC BY-SA. Présentation adaptée par Maths Explorer.}\par\medskip
\medskip

Une entreprise vend des smartphones d'un seul modèle \og haut de gamme \fg.

Le service marketing modélise le nombre de smartphones modèle \og haut de gamme\fg{}
vendus par trimestre en fonction du prix de vente $x$ par la fonction $N$ définie par 

\[N(x) = 100\e^{-2x}\,  \text{où :}\]


\setlength\parindent{1cm}
\begin{itemize}
\item $x$ est le prix de vente en \textbf{milliers d'euros} d'un smartphone modèle \og haut de gamme \fg.

 Le prix du smartphone modèle \og haut de gamme\fg{} est compris entre 400~\euro{} et \np{2000}~\euro{} ; on a donc $x \in [0,4~;~2]$.
\item $N(x)$ est le nombre de smartphones modèle \og haut de gamme\fg{} vendus trimestriellement en \textbf{millions d'unités}.
\end{itemize}
\setlength\parindent{0cm}

\smallskip

\begin{enumerate}
\item Si le service commercial fixe le prix de vente de ce smartphone modèle \og haut de gamme\fg{} à \np{1000}~\euro, quel sera le nombre de smartphones vendus trimestriellement ? 

On arrondira le résultat à mille unités.
\end{enumerate}

La recette trimestrielle $R(x)$ est obtenue en multipliant le nombre de smartphones modèle \og haut de gamme\fg{} vendus par le prix de vente. 

On obtient $R(x) = x \times  N(x)$ en \textbf{milliards d'euros}.

Le coût de production en milliards d'euros en fonction du nombre de smartphones modèle \og haut de gamme\fg{} fabriqués est modélisé par la fonction $C$ définie par $C(x) = 0,4 \times N(x)$ où $x$ est le prix de vente en \textbf{milliers d'euros}.

Le bénéfice est obtenu en calculant la différence entre la recette et le coût de production.

\smallskip

\begin{enumerate}[resume]
\item Vérifier que le bénéfice trimestriel peut être estimé à $8,120$ milliards d'euros pour un prix de vente \np{1000}~\euro.
\item Montrer que le bénéfice trimestriel s'exprime en milliards d'euros en fonction du prix de vente $x$ en milliers d'euros par : $B(x) = (100x - 40) \e^{-2x}.$
\item On admet que pour tout réel $x \in  [0,4~;~2],\: B'(x) = (180 - 200x )\e^{-2x}$.

 Étudier les variations de la fonction $B$ sur l'intervalle [0,4~;~2].
\item À quel prix faut-il vendre ces smartphones pour assurer un bénéfice maximal ?
\end{enumerate}

\end{document}
